\section{Related Work}
\label{sec:related}

\para{Instruction-following evaluation.}
The evaluation of LLM instruction following has received considerable attention. IFEval~\citep{zhou2023ifeval} introduced the concept of \emph{verifiable instructions}---constraints that can be checked automatically---and defined 25 instruction types covering keywords, length, format, and content constraints. FollowBench~\citep{jiang2023followbench} extended this with multi-level constraint difficulty, showing that models degrade as constraints become more complex. COLLIE~\citep{yao2023collie} proposed a grammar-based framework for compositional constraint specification at word, sentence, paragraph, and passage levels. \citet{qin2024benchmarking} studied how performance degrades with multiple constraints composed together. These benchmarks focus on \emph{imperative} instructions (``always do X'') and do not isolate the conditional reasoning required by if-then instructions. Our work specifically targets the conditional dimension.

\para{Sequential and compositional instruction following.}
The SIFo benchmark~\citep{chen2024sifo} evaluates sequential instruction following, where each step depends on the previous one. Its security-rules task---which requires following conditional access rules---is the closest existing proxy for conditional instruction evaluation. However, it covers only one narrow task type and does not systematically vary the type of conditional logic. CoDI-Eval~\citep{chen2024codieval} evaluates controllability under diversified constraint expressions, showing significant gaps between open-source and closed-source models. We build on these ideas by creating a taxonomy of 8 conditional instruction types and measuring performance across all of them.

\para{Internal mechanisms of instruction following.}
\citet{heo2024llms} discovered an ``instruction-following dimension'' in LLM embedding space that predicts compliance but, critically, does not generalize across instruction types. This finding directly motivates our hypothesis: if internal signals for instruction following are type-specific, then conditional instruction-following capacity should vary systematically across types. Our empirical results confirm this prediction.

\para{Improving instruction following.}
Several recent works propose methods to improve instruction adherence. \citet{kumar2025pseudocode} showed that expressing instructions as pseudo-code yields 8--21\% gains on instruction-following benchmarks, suggesting that structured representations help models parse conditional logic. \citet{wu2025thinking} found that chain-of-thought reasoning can sometimes \emph{hurt} instruction following, particularly for constraint adherence. \citet{zhang2024onlyif} demonstrated that generalization to unseen instructions requires diverse training data across semantic domains. These findings inform how practitioners might improve conditional instruction reliability beyond benchmark measurement.
